\MWChapter{Evidence and conformance status}{证据与符合性状态}
\label{ch:evaluation}

\section{How to read the status / 如何解读状态}

The central evidence is a proof-obligation manifest, not a benchmark suite. “Proved” means the
named obligation is bound to the recorded source and build evidence; “review-pending” means the
corresponding independent review and governance decision have not been recorded. The two labels
must appear together wherever the current core theory is summarized \citep{cantiluneClosure,cantiluneManifest}.

\section{Recorded theory evidence / 已记录的理论证据}

\begin{table}[H]
\centering
\caption{Current evidence statements and their deliberate boundaries.}
\label{tab:headline-results}
\begin{tabularx}{\textwidth}{@{}L{0.22\textwidth}L{0.23\textwidth}Y@{}}
\toprule
\rowcolor{MWSoft!70}
\textbf{Surface} & \textbf{Recorded evidence} & \textbf{Not established by that evidence} \\
\midrule
Generic core theory & 18/18 central obligations recorded \emph{proved} against the verified source and evidence chain & Independent QA-L4 review, FCP result, or ADR acceptance \\
P1a complete interface & A separate fixed-signature fourteen-event reference provides non-vacuity witnesses for supplied whole-LTS certificates & A certificate automatically generated for an arbitrary product rule family \\
P1c operation layer & Fifteen normative native event families and a sixty-operation registry have recorded strong late-\(\pi\) representatives & A public session API, provider integration, or runtime interoperability guarantee \\
D1-A boundary & Maximum-compatible D1-A closure records the selected lower-effect route and stated adequacy/full-abstraction scopes & A unified source-paper FMS model, unrestricted actual-Agent theorem, or reverse definability \\
Feedback layer & Finite-height closure, terminal classification, replay, and conditional hitting statements are part of the theory record & A production monitoring system \\
\bottomrule
\end{tabularx}
\end{table}

\section{Consistency is an event-level claim / 一致性是事件级声明}

For a recorded source change \(G\xrightarrow{e}G'\), the useful claim is not merely that four
dashboards can display related information. It is that every supplied projection certificate
relates the same event identity to a target step and preserves the declared observation boundary.
This is why the repository separates source events, target derivations, replay, terminal classes,
and product certificates \citep{cantiluneRFCTwo}.

\begin{equation}
  \operatorname{Replay}(G,e)=G'
  \quad\text{and}\quad
  \operatorname{Replay}\bigl(P_i(G),\Phi_i(e)\bigr)=P_i(G')
  \label{eq:replay}
\end{equation}

Equation \cref{eq:replay} is a target relation whose exact equality and certificate premises are
defined by the formal artifacts; it is not an assertion that every future implementation can
replay arbitrary external side effects.

\section{No empirical scorecard / 不提供经验指标}

The earlier runtime-template bars, latency curves, success rates, and safety-suite counts have
been removed from this report. They had no evidentiary relation to Cantilune's current formal
source tree. A future empirical report must name a released artifact, population, metrics,
baselines, raw evidence, and uncertainty before it can make performance or safety claims.

\begin{MWInsight}
Formal status and empirical status answer different questions. A kernel-checked theorem can
establish a stated mathematical consequence without establishing speed, usability, robustness,
or real-world safety; conversely, a benchmark cannot replace a missing semantic proof.
\end{MWInsight}
